\documentclass[11pt]{article}
\usepackage[margin=1in]{geometry}
\usepackage{amsmath,amssymb,mathtools,booktabs,array,microtype}
\usepackage[colorlinks=true,linkcolor=blue,citecolor=blue,urlcolor=blue]{hyperref}
\hypersetup{pdftitle={V1 neural-field coefficients: sources and derivation}}
\newcommand{\ee}{\mathrm e}
\newcommand{\Rea}{\operatorname{Re}}
\newcommand{\Ima}{\operatorname{Im}}
\newcommand{\diag}{\operatorname{diag}}
\setlength{\parindent}{0pt}
\setlength{\parskip}{5pt}
\begin{document}
\begin{center}\large\bfseries V1 neural-field coefficients: sources and derivation\end{center}
This note accompanies the V1 notebook and extends the model in \texttt{final\_report.pdf}, confirmed by the author as similar to the supplied
report. The gains below are reproducible phenomenological calibrations, not
estimates of physiological synaptic strengths.

\section{What is fixed by evidence, and what is assumed?}
Order the populations as $(E,PV,SST,VIP)$, with target rows and source columns.
The four-population sign/zero mask and the no-VIP restriction are
\begin{equation}
 \mathcal M_4=\begin{pmatrix}+&-&-&0\\+&-&-&0\\+&0&0&-\\+&0&-&0\end{pmatrix},
 \qquad \mathcal M_3=(\mathcal M_4)_{1:3,1:3}.
 \label{eq:mask}
\end{equation}
Every existing connection and its sign are preserved. Exact zeros represent
the chosen reduction, not proof that the corresponding biological pathways
never occur.

\textbf{Pfeffer et al.\ \cite{pfeffer}.} The measured inhibitory motif motivates
PV inhibition of E/PV, SST inhibition of the other retained populations, and
preferential VIP inhibition of SST. It does not fix our effective gain magnitudes.

\textbf{Billeh et al.\ \cite{billeh}.} Their Gaussian distance rule supplies the
scales $s_{E\to E}=114$, $s_{E\to PV}=92$, $s_{E\to SST}=103$,
$s_{PV\to E}=95$, $s_{SST\to E}=85$, and $s_{PV\to PV}\simeq120\,\mu$m.
Their substitutions for other inhibitory pairs and Htr3a connections motivate
our remaining widths; Htr3a is broader than VIP. Their feedforward model motivates
driving E/PV, although equal E/PV drive is our assumption.

\textbf{Campagnola et al.\ \cite{campagnola}.} Cell-class/layer dependence of
connection probability, strength and dynamics limits the interpretation of a
single coarse matrix. No numerical gain is copied from this study.

\textbf{Hahn et al.\ \cite{hahn}.} Their time-constant sensitivity analysis
(Fig.~5E--H) motivates testing frequency sensitivity. It does \emph{not} establish
our exact time constants. No cited experiment uniquely determines those values.

\subsection{From spatial measurements to field coefficients}
For a unit-mass Gaussian, one spatial unit is $100\,\mu$m, and
\begin{equation}
 P_{ij}(r)=C_{ij}\ee^{-r^2/s_{ij}^2},\qquad
 \Sigma_{ij}=\frac{(s_{ij}/100)^2}{2},\qquad
 \widehat W_{ij}(q;p)=J_{ij}(p)\ee^{-\Sigma_{ij}q^2/2}.
 \label{eq:width}
\end{equation}
Thus $J_{ij}$ is an integrated gain, not a probability or Gaussian peak height.
The disk-average convention in \cite{billeh}, with $R=75\,\mu$m, gives
\begin{equation}
 P_{{\rm rep},ij}=\frac{1}{\pi R^2}\int_0^R2\pi r C_{ij}\ee^{-r^2/s_{ij}^2}\,dr
 =\frac{C_{ij}s_{ij}^2}{R^2}(1-\ee^{-R^2/s_{ij}^2}).
\end{equation}
Consequently $C_{ij}=P_{{\rm rep},ij}R^2/[s_{ij}^2(1-\ee^{-R^2/s_{ij}^2})]$.
If source density and effective per-connection strength were known, a planar
reduction would give $J_{ij}\propto\rho_j w_{ij}\pi s_{ij}^2C_{ij}$, with further
activity-unit scaling. The legacy probability array does not supply these
missing factors. The 1D marginal of the unit-mass 2D kernel preserves $J_{ij}$.

\newpage
\section{Complete coefficient specification}
The source-nonlinearity convention is
\begin{equation}
 \tau_i\partial_tu_i=-u_i+\sum_jW_{ij}*f_j(u_j)+I_i,\qquad
 f_j(u)=\ell(g_j(u-\theta_j))-\ell(-g_j\theta_j),\quad
 \ell(z)=\frac{1}{1+\ee^{-z}}.
 \label{eq:model}
\end{equation}
Here $u$ is a signed deviation variable, time is in ms, and
$L=47.6023$ model units $=4.76023$ mm. Retain the legacy activation parameters
and effective time scales:
\begin{center}
\begin{tabular}{lrrr}\toprule
Population&$\tau$ (ms)&$g$&$\theta$\\\midrule
E&10&3.47682761&0.28483799\\
PV&12&7.34413013&0.12068148\\
SST&35&3.74650527&0.40137356\\
VIP&20&6&0.7\\\bottomrule
\end{tabular}
\end{center}
A one-pole filter has $h_i(t)=\tau_i^{-1}\ee^{-t/\tau_i}$ for $t\geq0$, so
$\int_0^\infty t h_i(t)dt=\tau_i$. No measured population impulse responses
were available to identify these moments. We therefore retain the ratios
$1:1.2:3.5:2$, with a 10 ms reference, as explicit hypotheses. They are not
measured membrane or synaptic-decay constants. Holding them and $g,\theta$ fixed
isolates the change in connectivity; their decimal precision is not experimental
precision.

The width matrix, including dynamically irrelevant placeholders on zero edges, is
\begin{equation}
 s=\begin{pmatrix}114&95&85&85\\92&120&120&120\\103&120&120&120\\103&120&120&120\end{pmatrix}\mu\mathrm m,
 \qquad \Sigma=(s/100)^2/2\quad\text{entrywise}.
\end{equation}
The $120\,\mu$m SST--VIP/SST--PV and $103\,\mu$m E--VIP scales are
extrapolations in the reduction, not VIP-specific measurements.

Set $J(p)=(0.5+p)G$, equivalently $A=0.5G$ and $A_p=G$. The calibrated gains are
\begin{align}
 G_3&=\begin{pmatrix}
 5.87511&-1.36667&-6.36306\\
 4.74411&-0.99825&-8.11851\\
 2.64504&0&0
 \end{pmatrix},\\
 G_4&=\begin{pmatrix}
 6.13537&-1.42384&-6.23105&0\\
 5.13746&-0.89269&-7.81025&0\\
 2.54421&0&0&-16.04588\\
 1.65365&0&-3.21101&0
 \end{pmatrix}.
 \label{eq:gains}
\end{align}
The notebook uses the full-precision values in
\path{simulations/output/v1_extension/calibration.json}. The printed matrices
are rounded. The no-VIP block is calibrated separately; comparing these matched
models is not an isolated causal VIP-ablation experiment.

\newpage
\section{Reproducible derivation of the new gains}
The legacy full-scale matrix is
\begin{equation}
 G_{\rm old}=\begin{pmatrix}
6.22009&-1.45455&-6.48944&0\\5.09798&-0.90615&-8.04066&0\\
2.67174&0&0&-16\\1.68844&0&-3.13584&0
\end{pmatrix}.
\end{equation}
Its original optimization history cannot be reconstructed. Instead, take this
matrix (or its principal $3\times3$ block) as a prior and solve
\begin{equation}
 \min_{-2\leq z_{ij}\leq2}\sum_{G_{{\rm old},ij}\ne0}z_{ij}^2,
 \qquad G_{ij}=G_{{\rm old},ij}\ee^{z_{ij}}\quad\text{on retained edges},
 \label{eq:optimization}
\end{equation}
with zeros fixed, subject to the four spectral targets
\begin{align}
 \Rea\lambda_d(q_0,p_0)&=0,&\partial_q\Rea\lambda_d(q_0,p_0)&=0,\notag\\
 \Ima\lambda_d(q_0,p_0)&=0.08,&\max\Rea\sigma(M(0,p_0))&=-0.001,\quad
 q_0=6\pi/L,\ p_0=0.25. \label{eq:targets}
\end{align}
The matrix $M$ is derived in Report 2. These are design targets: mode 3,
frequency $12.7324$ Hz and a damped homogeneous mode, not measured V1 values.
SLSQP starts at $z=0$, uses tolerance $10^{-12}$ and a central $q$-difference
of $10^{-4}$. It finds a local constrained optimum; uniqueness/global optimality
is not claimed. Numerical checks find the first crossing at $(p_0,k_0)=(0.25,3)$,
verify the square-lattice shells and a stable Gaussian tail, and obtain
negative real cubic saturation coefficients. Unretuned VIP deletion instead
gives a homogeneous Hopf at $p\simeq0.153726$.

\textbf{Forcing.} At $p=p_0-0.0005$, drive E/PV equally at $2\omega_0$ and/or
spatial mode 6. A real cosine of peak $d$ has complex harmonic $d/2$.
Solving for leading-order amplitude-Jacobian growth $10^{-4}$/ms gives
\begin{center}\small
\begin{tabular}{lrrr}\toprule
Model&Temporal peak&Spatial peak&Each combined peak\\\midrule
3 populations&0.0022181&0.0061721&0.0016479\\
4 populations&0.0021273&0.0028546&0.0012339\\\bottomrule
\end{tabular}\end{center}
Direct simulations check the resulting patterns separately from this asymptotic
prediction. With gains fixed, $\pm20\%$ changes in E/PV/SST time constants often
remove the finite-wave-first scenario: these are conditional existence examples,
not a robust physiological fit.

\begingroup\footnotesize
\begin{thebibliography}{4}\setlength{\itemsep}{0pt}
\bibitem{pfeffer} C.~K. Pfeffer et al. (2013), Inhibition of inhibition in visual cortex,
\emph{Nature Neuroscience} 16, 1068--1076. \href{https://doi.org/10.1038/nn.3446}{doi:10.1038/nn.3446}.
\bibitem{billeh} Y.~N. Billeh et al. (2020), Systematic Integration of Structural and
Functional Data into Multi-scale Models of Mouse Primary Visual Cortex,
\emph{Neuron} 106, 388--403.e18; Fig.~4, Methods e10--e11.
\href{https://doi.org/10.1016/j.neuron.2020.01.040}{doi:10.1016/j.neuron.2020.01.040}.
\bibitem{campagnola} L. Campagnola et al. (2022), Local connectivity and synaptic
dynamics in mouse and human neocortex, \emph{Science} 375, eabj5861.
\href{https://doi.org/10.1126/science.abj5861}{doi:10.1126/science.abj5861}.
\bibitem{hahn} G. Hahn et al. (2022), Rate and oscillatory switching dynamics of a
multilayer visual microcircuit model, \emph{eLife} 11, e77594; Fig.~5E--H.
\href{https://doi.org/10.7554/eLife.77594}{doi:10.7554/eLife.77594}.
\end{thebibliography}\endgroup
\end{document}
